\documentclass[../../../include/open-logic-section]{subfiles}

\begin{document}

\olfileid{sth}{card-arithmetic}{ch}

\olsection{Hipotesis Kontinum}

Asil sadurungé mènèhi pratandha, lan pancèn bener, yèn pemangkatan
kardinal bakal cukup \emph{gampang} saupama kardinal tanpa wates mesthi
sesambungan kanthi prasaja karo pangkat $2$, yaiku, miturut
\olref[opps]{lem:SizePowerset2Exp}, karo operasi njupuk himpunan pangkat.
Nanging kita ora bisa nganggep yèn kardinal tanpa wates \emph{pancèn} tansah
sesambungan kanthi prasaja karo himpunan pangkat.

Kanggo miwiti njlentrehaké prakara iki, kita ngenalaké notasi kang migunani.

\begin{defn}
Yèn $\cardsucc{\cardfont{a}}$ iku kardinal paling cilik kang luwih gedhé
kanthi ketat tinimbang $\cardfont{a}$, kita netepaké rong runtutan
transfinit kaya mangkéné:
\begin{align*}
	\aleph_{0} &\defis \omega & 		
	\beth_{0} &\defis \omega\\
	\aleph_{\alpha \ordplus 1} &\defis \cardsucc{(\aleph_{\alpha})} &
	\beth_{\alpha+1} &\defis \cardexpo{2}{\beth_{\alpha}}\\
	\aleph_{\alpha} &\defis \bigcup_{\beta< \alpha} \aleph_{\beta} &
	\beth_{\alpha} &\defis \bigcup_{\beta < \alpha}\beth_{\beta} & \text{nalika $\alpha$ ordinal limit}.
	\end{align*}
\end{defn}

Dhéfinisi $\cardsucc{\cardfont{a}}$ iki sah: miturut
\olref[cardinals][classing]{lem:NoLargestCardinal}, kanggo saben
kardinal $\cardfont{a}$ ana kardinal kang luwih gedhé tinimbang
$\cardfont{a}$, lan Induksi Transfinit njamin anané kardinal kang
\emph{paling cilik} ing antarané kang luwih gedhé tinimbang
$\cardfont{a}$.
Dhéfinisi pérangan liya saka runtutan $\aleph$ lan $\beth$ diwènèhaké
lumantar rekursi transfinit.

Cantor ngenalaké notasi ``$\aleph$''; iki diwaca \emph{alef}, aksara
kapisan ing abjad Ibrani lan aksara kapisan ing tembung Ibrani kanggo
``tanpa wates''. Peirce ngenalaké notasi ``$\beth$''; iki diwaca
\emph{bet}, aksara kapindho ing abjad Ibrani.\footnote{Peirce migunakaké
notasi iki ing layang marang Cantor ing Desember 1900. Sayangé, ing
layang iku Peirce uga mènèhi panalaran kang klèru yèn $\beth_\alpha$
ora ana kanggo $\alpha \geq \omega$.} Notasi-notasi iki mènèhi kita
kardinal tanpa wates.

\begin{prop}
$\aleph_\alpha$ lan $\beth_\alpha$ iku kardinal kanggo saben
ordinal $\alpha$.
\end{prop}

\begin{proof}
Kaloro pratelan iki kabukten kanthi induksi transfinit kang prasaja.
$\aleph_0 = \beth_0 = \omega$ iku kardinal miturut
\olref[cardinals][classing]{omegaisacardinal}. Yèn $\aleph_\alpha$ lan
$\beth_\alpha$ padha kardinal, mula $\aleph_{\alpha+1}$ lan
$\beth_{\alpha+1}$ kanthi gamblang didhéfinisikaké minangka kardinal.
Ing tahap limit, gabungan sawijiné himpunan kardinal uga kardinal,
miturut \olref[cardinals][classing]{unioncardinalscardinal}.
\end{proof}
\noindent
Kajaba iku, saben kardinal tanpa wates iku sawijining $\aleph$:

\begin{prop}
Yèn $\cardfont{a}$ kardinal tanpa wates, mula $\cardfont{a} =
\aleph_\gamma$ kanggo sawijining $\gamma$ kang tunggal.
\end{prop}

\begin{proof}
Kita migunakaké induksi transfinit marang kardinal tanpa wates.
Kasus wiwitané $\cardfont{a}=\omega=\aleph_0$. Kanggo langkah
induksi, anggepen yèn saben $\omega \leq \cardfont{b} < \cardfont{a}$
wis nduwèni wujud $\cardfont{b} = \aleph_{\gamma_\cardfont{b}}$.
Yèn $\cardfont{a} = \cardsucc{\cardfont{b}}$ kanggo sawijining
$\cardfont{b}$ tanpa wates, mula $\cardfont{a} =
\cardsucc{(\aleph_{\gamma_\cardfont{b}})}=
\aleph_{\gamma_\cardfont{b}+1}$. Yèn $\cardfont{a}$ dudu penerus
sawijining kardinal, mula amarga kardinal uga ordinal,
$\cardfont{a} = \bigcup_{\omega \leq \cardfont{b} < \cardfont{a}}
\cardfont{b} =
\bigcup_{\omega \leq \cardfont{b} < \cardfont{a}}{\aleph_{\gamma_\cardfont{b}}}$.
Dadi $\cardfont{a} = \aleph_\gamma$ kanthi $\gamma =
\bigcup_{\omega \leq \cardfont{b} < \cardfont{a}}\gamma_\cardfont{b}$.
Indeksé tunggal amarga runtutan $\aleph$ mundhak kanthi ketat.
\end{proof}

Amarga saben kardinal tanpa wates iku sawijining $\aleph$, kita banjur
bisa takon: apa saben kardinal tanpa wates uga sawijining~$\beth$?
Yèn pancèn \emph{mangkono}, kardinal tanpa wates mesthi sesambungan kanthi
prasaja karo operasi njupuk himpunan pangkat. Pancen, kita bakal
nduwèni pratelan iki:

\begin{defish}
\emph{Hipotesis Kontinum Umum} (GCH). $\aleph_\alpha  = \beth_\alpha$, kanggo saben $\alpha$.
\end{defish}

Kajaba iku, yèn GCH bener, kita bisa nyederhanakaké pemangkatan kardinal
kanthi cukup adoh. Mliginé, kanggo $\cardfont{a}$ tanpa wates lan
$0<\cardfont{b}<\cardfont{a}$, kita bisa mbuktekaké yèn
$\cardexpo{\cardfont{a}}{\cardfont{b}}$ kaapit déning
$\cardfont{a}\leq \cardexpo{\cardfont{a}}{\cardfont{b}} \leq
\cardsucc{\cardfont{a}}$. Sabanjuré, kita bisa mènèhi syarat-syarat
pas kang nemtokaké endi saka rong kamungkinan kang kelakon: apa
$\cardfont{a} = \cardexpo{\cardfont{a}}{\cardfont{b}}$ utawa
$\cardexpo{\cardfont{a}}{\cardfont{b}} =
\cardsucc{\cardfont{a}}$.\footnote{Syaraté ditemtokaké déning
\emph{kofinalitas}.}

Nanging GCH iku sawijining \emph{hipotesis}, dudu \emph{téoréma}.
\citet{Godel1938} mbuktekaké yèn $\ZFC$ konsisten, mula
$\ZFC + \text{GCH}$ uga konsisten. Nanging banjur kabukten yèn kita uga bisa
nambahaké $\lnot$GCH marang $\ZFC$. Gatekna \emph{kasus} GCH kang paling
prasaja nanging ora sepele iki:

\begin{defish}
\emph{Hipotesis Kontinum} (CH). $\aleph_1 = \beth_1$.
\end{defish}

\citet{Cohen1963} mbuktekaké yèn $\ZFC$ konsisten, mula
$\ZFC + \lnot\text{CH}$ uga konsisten. Mula Hipotesis Kontinum independen
saka $\ZFC$.

Jeneng Hipotesis Kontinum asalé saka tembung ``kontinum'', kang uga
nyebut garis bilangan riil, $\Real$.
\olref[opps]{continuumis2aleph0} nyatakaké yèn $\card{\Real} =
\beth_1$. Dadi Hipotesis Kontinum nyatakaké yèn ora ana kardinal
ing antarané kardinalitas bilangan asli, $\aleph_0 = \beth_0$,
lan kardinalitas kontinum, $\beth_1$.

Amarga (G)CH \emph{independen} saka $\ZFC$, apa kang kudu kita
wangsuli babagan \emph{bebeneré}? Bab iki bisa dirembug dawa,
lan \emph{akèh} panaliten teori himpunan kang subur bubar iki nyinaoni
prakara kasebut. Nanging ana rong bab cekak kang perlu digatèkaké.

Kapisan, saka asil independensi formal iki ora \emph{langsung}
katut yèn nilai bebeneré GCH utawa CH \emph{ora bisa ditemtokaké}.
Mbokmenawa kita mung perlu nambahaké aksioma kang katon lumrah
lan bisa netepaké jawabané. G\"odel dhéwé nyaranaké tanggapan kaya iki.

Kapindho, independensi CH saka $\ZFC$ pancèn
\emph{nggumunaké}, nanging ora \emph{ora bisa dipercaya} ing teges
harfiah. Pokoké, saka apa kang dikandhakake $\ZFC$ waé,
lumaku saka sawijining kardinal menyang penerusé bisa mbutuhaké
piranti kang luwih alus tinimbang mung njupuk himpunan pangkat.
 %Operasi njupuk himpunan pangkat mindhah kita saka siji tahap
 %hirarki menyang tahap penerusé (saka $V_{\alpha}$ menyang
 %$V_{\alpha+1}$).

Sawisé rong pengamatan mau, yèn panjenengan kepengin ngerti luwih
akèh, panjenengan kudu maca panemu para filsuf lan matématikawan kang
melu rembugan iki.\footnote{Minangka wiwitan, panjenengan bisa maca
\citet[\S15.6]{Potter2004}.}

\end{document}
